% GENERATED FILE. Edit canonical lessons, then run npm run generate:books.
% canonical-source-hash: fnv1a64:630b5a97c2c82ac1
% canonical-lessons: BN-C470-tak, BN-C470-khushki, BN-C470-ukun, BN-C470-jhora, BN-C470-khopa

\chapter{Hair and Scalp}
\label{ch:bn-hair-and-scalp}
\hlchaptermodality{470}

\begin{chapteropening}
\textbf{By the end of this chapter:} \emph{I can talk about baldness, dandruff, lice and greying hair.}

Complete the last lesson of chapter 470: I can talk about baldness, dandruff, lice and greying hair.
\end{chapteropening}

\section[\texorpdfstring{\bn{টাক} (ṭāk)}{টাক (ṭāk)}]{\bn{টাক} (ṭāk) — a bald patch; baldness}

\label{lesson:BN-C470-tak}



\begin{quote}
Before the new one: say the Bengali for a stain, then the Bengali for body hair.
\end{quote}

\subsection*{You'll want to know: \bn{টাক}}
\textbf{\bn{টাক}} — \emph{ṭāk} — \textquotedblleft{}a bald patch; baldness\textquotedblright{}.

\begin{quote}
\bn{চল্লিশ} \bn{বছরে} \bn{বাবার} \bn{মাথায়} \bn{টাক} \bn{পড়ল।}
\end{quote}

\begin{quote}
\emph{chôllish bôchhôre bābār māthāy ṭāk pôṛlô.}
\end{quote}

\begin{quote}
At forty, my father started going bald.
\end{quote}

\bn{টাক} \bn{পড়া}, to go bald.

\begin{cousinweb}
\textbf{\bn{টাক}} (\emph{ṭāk}) is listed by CDIAL for Bengali, as a bald patch, under a reconstructed \emph{ṭakka}, \textquotedblleft{}bald\textquotedblright{} (CDIAL 5422); it is not the \bn{টাকা} (taught) of money, which is another entry.
\end{cousinweb}

\begin{grammarlens}[title={What you've built}]
One more word for describing what you lived and felt.
\end{grammarlens}

\subsection*{Your turn}
\begin{itemize}
\raggedright
  \item \emph{Say it:} \emph{ṭāk}
  \item \emph{Say it:} \emph{ṭāk}, once more
  \item \emph{Say it:} \emph{chôllish bôchhôre bābār māthāy ṭāk pôṛlô.} — At forty, my father started going bald.
  \item \emph{Recall:} say the Bengali for a stain, then the Bengali for body hair, then say \emph{ṭāk} again
\end{itemize}

\begin{tcolorbox}[breakable,colback=teal!4,colframe=teal!35!black,title={Before you move on}]
What does \bn{টাক} mean? (\textquotedblleft{}A bald patch; baldness\textquotedblright{}.) Say it once more.
\end{tcolorbox}
\section[\texorpdfstring{\bn{খুশকি} (khushki)}{খুশকি (khushki)}]{\bn{খুশকি} (khushki) — dandruff}

\label{lesson:BN-C470-khushki}



\begin{quote}
Before the new one: say the Bengali for body hair, then the Bengali for baldness.
\end{quote}

\subsection*{You'll want to know: \bn{খুশকি}}
\textbf{\bn{খুশকি}} — \emph{khushki} — \textquotedblleft{}dandruff\textquotedblright{}.

\begin{quote}
\bn{শীতে} \bn{অরুণের} \bn{মাথায়} \bn{খুশকি} \bn{হয়।}
\end{quote}

\begin{quote}
\emph{shīte ôruner māthāy khushki hôy.}
\end{quote}

\begin{quote}
In winter Arun gets dandruff.
\end{quote}

\begin{cousinweb}
\textbf{\bn{খুশকি}} (\emph{khushki}) is usually traced to Persian \emph{khushkī}, \textquotedblleft{}dryness\textquotedblright{}.
\end{cousinweb}

\begin{grammarlens}[title={What you've built}]
One more word for describing what you lived and felt.
\end{grammarlens}

\subsection*{Your turn}
\begin{itemize}
\raggedright
  \item \emph{Say it:} \emph{khushki}
  \item \emph{Say it:} \emph{khushki}, once more
  \item \emph{Say it:} \emph{shīte ôruner māthāy khushki hôy.} — In winter Arun gets dandruff.
  \item \emph{Recall:} say the Bengali for body hair, then the Bengali for baldness, then say \emph{khushki} again
\end{itemize}

\begin{tcolorbox}[breakable,colback=teal!4,colframe=teal!35!black,title={Before you move on}]
What does \bn{খুশকি} mean? (\textquotedblleft{}Dandruff\textquotedblright{}.) Say it once more.
\end{tcolorbox}
\section[\texorpdfstring{\bn{উকুন} (ukun)}{উকুন (ukun)}]{\bn{উকুন} (ukun) — a louse; lice}

\label{lesson:BN-C470-ukun}



\begin{quote}
Before the new one: say the Bengali for baldness, then the Bengali for dandruff.
\end{quote}

\subsection*{You'll want to know: \bn{উকুন}}
\textbf{\bn{উকুন}} — \emph{ukun} — \textquotedblleft{}a louse; lice\textquotedblright{}.

\begin{quote}
\bn{স্কুল} \bn{থেকে} \bn{ফিরে} \bn{মেয়েটার} \bn{মাথায়} \bn{উকুন} \bn{হল।}
\end{quote}

\begin{quote}
\emph{skul theke phire meyeṭār māthāy ukun hôlô.}
\end{quote}

\begin{quote}
After coming home from school, the girl had lice.
\end{quote}

\begin{cousinweb}
\textbf{\bn{উকুন}} (\emph{ukun}) comes from Sanskrit \emph{utkuṇa}, \textquotedblleft{}a louse\textquotedblright{} (CDIAL 1728), which CDIAL lists in this Bengali form.
\end{cousinweb}

\begin{grammarlens}[title={What you've built}]
One more word for describing what you lived and felt.
\end{grammarlens}

\subsection*{Your turn}
\begin{itemize}
\raggedright
  \item \emph{Say it:} \emph{ukun}
  \item \emph{Say it:} \emph{ukun}, once more
  \item \emph{Say it:} \emph{skul theke phire meyeṭār māthāy ukun hôlô.} — After coming home from school, the girl had lice.
  \item \emph{Recall:} say the Bengali for baldness, then the Bengali for dandruff, then say \emph{ukun} again
\end{itemize}

\begin{tcolorbox}[breakable,colback=teal!4,colframe=teal!35!black,title={Before you move on}]
What does \bn{উকুন} mean? (\textquotedblleft{}A louse; lice\textquotedblright{}.) Say it once more.
\end{tcolorbox}
\section[\texorpdfstring{\bn{ঝরা} (jhôrā)}{ঝরা (jhôrā)}]{\bn{ঝরা} (jhôrā) — to fall out (of hair); to fall away; to drip}

\label{lesson:BN-C470-jhora}



\begin{quote}
Before the new one: say the Bengali for dandruff, then the Bengali for a louse.
\end{quote}

\subsection*{You'll want to know: \bn{ঝরা}}
\textbf{\bn{ঝরা}} — \emph{jhôrā} — \textquotedblleft{}to fall out (of hair); to fall away; to drip\textquotedblright{}.

\begin{quote}
\bn{চিন্তায়} \bn{অরুণের} \bn{চুল} \bn{ঝরছে।}
\end{quote}

\begin{quote}
\emph{chintāy ôruner chul jhôrchhe.}
\end{quote}

\begin{quote}
Arun's hair is falling out from worry.
\end{quote}

Past (he/she …): \bn{সেই} \bn{গাছের} \bn{পাতা} \bn{ঝরল}, \emph{sei gāchher pātā jhôrlô}.

\begin{cousinweb}
\textbf{\bn{ঝরা}} (\emph{jhôrā}) is listed by CDIAL for Bengali, as oozing, dripping and leaves falling, under a reconstructed form (CDIAL 5346); \bn{ঝরনা} (taught), a waterfall, is probably its relative.
\end{cousinweb}

\begin{grammarlens}[title={What you've built}]
One more word for describing what you lived and felt.
\end{grammarlens}

\subsection*{Your turn}
\begin{itemize}
\raggedright
  \item \emph{Say it:} \emph{jhôrā}
  \item \emph{Say it:} \emph{jhôrā}, once more
  \item \emph{Say it:} \emph{chintāy ôruner chul jhôrchhe.} — Arun's hair is falling out from worry.
  \item \emph{Recall:} say the Bengali for dandruff, then the Bengali for a louse, then say \emph{jhôrā} again
\end{itemize}

\begin{tcolorbox}[breakable,colback=teal!4,colframe=teal!35!black,title={Before you move on}]
What does \bn{ঝরা} mean? (\textquotedblleft{}To fall out (of hair); to fall away; to drip\textquotedblright{}.) Say it once more.
\end{tcolorbox}
\section[\texorpdfstring{\bn{খোঁপা} (khõpā)}{খোঁপা (khõpā)}]{\bn{খোঁপা} (khõpā) — a bun; a knot of hair}

\label{lesson:BN-C470-khopa}



\begin{quote}
Before the new one: say the Bengali for a louse, then the Bengali for to fall out.
\end{quote}

\subsection*{You'll want to know: \bn{খোঁপা}}
\textbf{\bn{খোঁপা}} — \emph{khõpā} — \textquotedblleft{}a bun; a knot of hair\textquotedblright{}.

\begin{quote}
\bn{মীরা} \bn{চুলে} \bn{একটা} \bn{খোঁপা} \bn{বাঁধল।}
\end{quote}

\begin{quote}
\emph{mīrā chule ekṭā khõpā bā̃dhlô.}
\end{quote}

\begin{quote}
Mira put her hair up in a bun.
\end{quote}

\bn{খোঁপা} \bn{বাঁধা}, to put one's hair up in a bun.

\begin{cousinweb}
\textbf{\bn{খোঁপা}} (\emph{khõpā}) is listed by CDIAL for Bengali, as a knot of hair, under a reconstructed \emph{khōppa}, \textquotedblleft{}a bunch of hair\textquotedblright{} (CDIAL 3939), which CDIAL takes to be a borrowing from Dravidian.
\end{cousinweb}

\begin{grammarlens}[title={What you've built}]
One more word for describing what you lived and felt.
\end{grammarlens}

\subsection*{Your turn}
\begin{itemize}
\raggedright
  \item \emph{Say it:} \emph{khõpā}
  \item \emph{Say it:} \emph{khõpā}, once more
  \item \emph{Say it:} \emph{mīrā chule ekṭā khõpā bā̃dhlô.} — Mira put her hair up in a bun.
  \item \emph{Recall:} say the Bengali for a louse, then the Bengali for to fall out, then say \emph{khõpā} again
\end{itemize}

\begin{tcolorbox}[breakable,colback=teal!4,colframe=teal!35!black,title={Before you move on}]
What does \bn{খোঁপা} mean? (\textquotedblleft{}A bun; a knot of hair\textquotedblright{}.) Say it once more.
\end{tcolorbox}
